\chapter{Conclusion}

In this thesis, we have explored the classification of variable stars using \linebreak transformer-based models.
Theoretical foundations concerning both the astronomical and machine learning sides were established in the initial chapters, the astronomical knowledge of the problem was later crucial for designing a~suitable representation of the data for the transformer models.

Then, a separate experiment concerning classification of ECG signals was conducted to gain practical experience with transformer-based classifiers \linebreak on~a~well-documented task -- the MIT BIH dataset.
The experiment was successful and the results were competitive when compared with both transformer and non-transformer based methods implemented by other authors on the same dataset.

For light-curve classification, we used a manually labeled dataset derived from observations from the TESS mission.
We created a robust preprocessing pipeline that extracted multiple views of different resolutions, allowing the model to capture patterns from different timescales.
Additionally, each light curve was also split into multiple segments which increased the dataset size and improved training performance.

Next, we implemented a modular model architecture that could effectively handle the input data with multiple channels.
Extensive experimentation with different model configurations was performed on both transformers trained from scratch and fine-tunes of pretrained models.
The best performing model was a fine-tune of the Chronos 2 pretrained TSFM and has shown promising results for further applications.

\section{Viability of Results}

The thesis is a part of an ongoing research by the Astronomical Institute of Czech Academy of Sciences, whose purpose is the implementation of an~automated classifier that can help identify interesting objects, especially exoplanets.

Existing pipelines for the identification of exoplanets produce large amounts of false positives, usually misclassifying a variable star as an exoplanet.
Additional filtering of these results is therefore necessary and the proposed model (Task B) could further reduce the number of false positives by correctly identifying pulsating and rotating variables (exoplanets are extremely rare and their light curves are very similar to those of eclipsing binaries).

Furthermore, two important missions regarding light curves are scheduled in the near future -- Gaia mission DR4 \parencite{Prusti2016} (new data release scheduled for December 2026) and PLATO mission \parencite{Rauer2025} (spacecraft launch scheduled for January 2027).
These missions will produce large amounts of light curves that will need to be processed and our models could either directly help the filtering or serve as a foundation for more research.

\section{Further Improvements}

Advanced model architectures like gated FFNs, FFT-based attention and hybrid fusion techniques are used in modern transformer architectures and could be utilized to improve the model's performance.
Unsupervised pretraining methods using unlabeled light curves are also a promising option that could help the model learn the general structure of the data without the need for limited labeled sets.

Furthermore, having a larger labeled dataset for training would also improve the predictive performance.
